\documentclass[11pt, a4paper]{article}

% ── Codificación y Lenguaje ──────────────────────────────────────────────────
\usepackage[utf8]{inputenc}
\usepackage[T1]{fontenc}
\usepackage[spanish, es-tabla]{babel}

% ── Geometría ───────────────────────────────────────────────────────────────
\usepackage[top=2.5cm, bottom=2.5cm, left=2.5cm, right=2.5cm, headheight=24pt]{geometry}

% ── Matemáticas ─────────────────────────────────────────────────────────────
\usepackage{amsmath, amssymb, amsthm}
\usepackage{mathtools}

% ── Colores y Cajas ─────────────────────────────────────────────────────────
\usepackage{xcolor}
\usepackage[most]{tcolorbox}
\tcbuselibrary{skins, breakable, theorems}

% ── Tablas ──────────────────────────────────────────────────────────────────
\usepackage{booktabs}
\usepackage{array}
\usepackage{multirow}
\usepackage{longtable}

% ── Listas ──────────────────────────────────────────────────────────────────
\usepackage{enumitem}

% ── Encabezado y Pie de Página ──────────────────────────────────────────────
\usepackage{fancyhdr}

% ── Secciones ───────────────────────────────────────────────────────────────
\usepackage{titlesec}

% ── Gráficos y Diagramas ────────────────────────────────────────────────────
\usepackage{tikz}
\usetikzlibrary{shapes.geometric, arrows.meta, positioning, calc}
\usepackage{graphicx}

% ── Columnas múltiples ──────────────────────────────────────────────────────
\usepackage{multicol}

% ── Espaciado ───────────────────────────────────────────────────────────────
\usepackage{setspace}
\setlength{\parskip}{4pt}
\setlength{\parindent}{0pt}

% ── Paleta GOP ──────────────────────────────────────────────────────────────
\definecolor{gopblue}{RGB}{31, 78, 121}
\definecolor{gopcyan}{RGB}{70, 130, 180}
\definecolor{gopgreen}{RGB}{0, 100, 0}
\definecolor{gopgreenlight}{RGB}{230, 255, 230}
\definecolor{goporange}{RGB}{200, 100, 0}
\definecolor{goporangelight}{RGB}{255, 245, 210}
\definecolor{gopred}{RGB}{160, 0, 0}
\definecolor{gopredlight}{RGB}{255, 230, 230}
\definecolor{gopgray}{RGB}{90, 90, 90}
\definecolor{gopgraylight}{RGB}{245, 245, 240}
\definecolor{gopbluedef}{RGB}{0, 70, 140}
\definecolor{gopbluedeflight}{RGB}{220, 235, 255}

% ── Hipervínculos y TOC ─────────────────────────────────────────────────────
\usepackage[colorlinks=true, linkcolor=gopblue, urlcolor=gopblue]{hyperref}

% ── Entornos tcolorbox ───────────────────────────────────────────────────────
\newtcolorbox{definicion}[1]{
  enhanced, breakable,
  colback=gopbluedeflight, colframe=gopbluedef,
  fonttitle=\bfseries\small, title={Definición: #1}, coltitle=white,
  attach boxed title to top left={yshift=-2mm, xshift=6pt},
  boxed title style={colback=gopbluedef, colframe=gopbluedef, sharp corners},
  sharp corners=south, top=4mm, before skip=6pt, after skip=6pt,
}
\newtcolorbox{formulas}[1]{
  enhanced, breakable,
  colback=gopgreenlight, colframe=gopgreen,
  fonttitle=\bfseries\small, title={Fórmulas: #1}, coltitle=white,
  attach boxed title to top left={yshift=-2mm, xshift=6pt},
  boxed title style={colback=gopgreen, colframe=gopgreen, sharp corners},
  sharp corners=south, top=4mm, before skip=6pt, after skip=6pt,
}
\newtcolorbox{tipprueba}[1]{
  enhanced, breakable,
  colback=goporangelight, colframe=goporange,
  fonttitle=\bfseries\small, title={TIP PRUEBA: #1}, coltitle=white,
  attach boxed title to top left={yshift=-2mm, xshift=6pt},
  boxed title style={colback=goporange, colframe=goporange, sharp corners},
  sharp corners=south, top=4mm, before skip=6pt, after skip=6pt,
}
\newtcolorbox{trampa}[1]{
  enhanced, breakable,
  colback=gopredlight, colframe=gopred,
  fonttitle=\bfseries\small, title={TRAMPA/ADVERTENCIA: #1}, coltitle=white,
  attach boxed title to top left={yshift=-2mm, xshift=6pt},
  boxed title style={colback=gopred, colframe=gopred, sharp corners},
  sharp corners=south, top=4mm, before skip=6pt, after skip=6pt,
}
\newtcolorbox{ejercicio}[1]{
  enhanced, breakable,
  colback=gopgraylight, colframe=gopgray,
  fonttitle=\bfseries\small\color{black}, title={Ejercicio: #1},
  attach boxed title to top left={yshift=-2mm, xshift=6pt},
  boxed title style={colback=gopgray!40, colframe=gopgray, sharp corners},
  sharp corners=south, top=4mm, before skip=6pt, after skip=6pt,
}

% ── Encabezado y pie ─────────────────────────────────────────────────────────
\pagestyle{fancy}
\fancyhf{}
\fancyhead[L]{\small\textbf{GOP ICS3213} --- Pauta Interrogación 2 (1er Semestre 2020)}
\fancyhead[C]{}
\fancyhead[R]{\small\thepage}
\fancyfoot[L]{\footnotesize Solución Oficial --- Transcripción en LaTeX}
\fancyfoot[R]{\small\thepage}
\renewcommand{\headrulewidth}{0.4pt}
\renewcommand{\footrulewidth}{0.4pt}

% ── Estilo de secciones ──────────────────────────────────────────────────────
\titleformat{\section}
  {\color{gopblue}\Large\bfseries}{\color{gopblue}\thesection.}{0.5em}{}
  [\color{gopblue}\titlerule]
\titleformat{\subsection}
  {\color{gopblue}\normalsize\bfseries}{\color{gopblue}\thesubsection.}{0.5em}{}
\titleformat{\subsubsection}
  {\color{gopblue}\small\bfseries}{\color{gopblue}\thesubsubsection.}{0.5em}{}

\begin{document}

% ════════════════════════════════════════════════════════════════════════════
% PORTADA
% ════════════════════════════════════════════════════════════════════════════
\begin{titlepage}
  \centering
  \vspace*{2cm}
  {\Huge\bfseries\color{gopcyan} Solución Interrogación 2\par}
  \vspace{0.4cm}
  {\LARGE\bfseries Gestión de Operaciones\par}
  \vspace{0.3cm}
  {\large ICS3213 --- PUC Chile --- 1er Semestre 2020\par}
  \vspace{1.2cm}
  \rule{\textwidth}{0.4pt}
  \vspace{0.4cm}

  {\small Profesores:\par}
  {\large\bfseries Martín García \quad Alejandro Mac Cawley\par}
  \vspace{0.3cm}
  {\small Autores de Referencia de la Pauta:\par}
  {\small\textbf{César Meneses \quad Fabián Ortega}\par}
  \vspace{0.4cm}
  \rule{\textwidth}{0.4pt}
  \vspace{1cm}

  \vfill
\end{titlepage}

\newpage

% ════════════════════════════════════════════════════════════════════════════
% PARTE I: PREGUNTAS CORTAS DE EJERCICIO
% ════════════════════════════════════════════════════════════════════════════
\section{Parte I: Preguntas Cortas de Ejercicio (30 Puntos)}

\subsection*{Pregunta I: Modelo de Programación Matemática para MRP (15 Puntos)}

  \begin{tipprueba}{¿Cuándo usar este enfoque?}
    \textbf{Identificación:} Se solicita formular un modelo de optimización (MILP) para planificar la producción bajo una estructura MRP (lista de materiales BOM, capacidades de producción, inventarios iniciales y finales, y lead times). Además, hay costos de producción, mantención y condiciones lógicas como lotes mínimos de producción.\\
    \textbf{Qué aplicar:} Definir conjuntos, índices, variables continuas de inventario, lanzamientos y requerimientos, variables binarias para la activación de producción en nodos específicos, y formular restricciones de balance, explosión de la BOM y límites de capacidad.
  \end{tipprueba}

  \textbf{Enunciado:}
  Usted dispone de la demanda diaria del producto que vende la empresa, dada por $F_t$ para $t$ periodos. La receta del producto requiere $j$ productos o subproductos y está dada por el \textit{Bill of Material} (BOM). El BOM le indica que para cada producto o subproducto $j$, usted requiere $R_{j,k}$ unidades de subproducto o insumo $k$ ($(j, k) \in \text{BOM}$). 

  Para cada producto/subproducto/insumo $j$ usted dispone de $II_j$ unidades en inventario inicial, el lead time de producción o del proveedor es de $L_j$ días para cada $j$ y finalmente, cada estación de producción o proveedor tiene una capacidad máxima de producción o entrega diaria de $C_j$ unidades. El costo de producción de cada unidad es de $CP_j$ pesos por unidad y el costo de mantener inventario es de $CI_j$ pesos por día por unidad $j$. 

  Finalmente, debido a condiciones técnicas hay dos partes $x, z \in \text{BOM}$ que tienen niveles mínimos de producción $CMin_x$ y $CMin_z$ cuando se inician sus procesos productivos. Con esta información elabore un modelo de programación matemática que permita determinar el programa de MRP para cumplir con la demanda, cantidad de insumos o subproductos $k$ en cada periodo y también terminar con $IF_j$ unidades de inventario final de cada unidad $j$.

  \medskip
  \begin{ejercicio}{Solución}
  \textbf{Desarrollo de la Solución:}

  \subsection*{1. Conjuntos e Índices}
  \begin{itemize}
    \item $j, k \in \mathcal{J}$: Conjunto de todos los ítems (productos, subproductos e insumos). El producto final está denotado por el índice $1$.
    \item $t \in \{1, \dots, T\}$: Períodos de planificación (días).
    \item $(j, k) \in \text{BOM}$: Relación que indica que el ítem $j$ requiere del ítem $k$ como componente directo.
  \end{itemize}

  \subsection*{2. Parámetros}
  \begin{itemize}
    \item $F_t$: Demanda externa del producto final ($1$) en el día $t$.
    \item $R_{j,k}$: Cantidad de unidades del componente $k$ requeridas para producir $1$ unidad del ítem $j$.
    \item $II_j$: Inventario inicial del ítem $j$ en el período $0$.
    \item $IF_j$: Inventario final requerido del ítem $j$ al final del período $T$.
    \item $L_j$: \textit{Lead Time} (tiempo de espera) de producción o compra del ítem $j$.
    \item $C_j$: Capacidad máxima de producción o entrega diaria del ítem $j$.
    \item $CMin_x, CMin_z$: Producción mínima requerida de los componentes $x$ y $z$ si se inicia su producción.
    \item $CP_j$: Costo unitario de producción del ítem $j$.
    \item $CI_j$: Costo diario de mantener $1$ unidad del ítem $j$ en inventario.
    \item $M$: Un número real lo suficientemente grande (\textit{Big-M}).
  \end{itemize}

  \subsection*{3. Variables de Decisión}
  \begin{itemize}
    \item $GR_{j,t} \ge 0$: Requerimientos brutos (\textit{Gross Requirements}) del ítem $j$ en el período $t$.
    \item $POR_{j,t} \ge 0$: Lanzamiento de orden planificada (\textit{Planned Order Release}) del ítem $j$ en el período $t$.
    \item $I_{j,t} \ge 0$: Inventario disponible del ítem $j$ al final del período $t$.
    \item $B_{i,t} \in \{0, 1\}$: Variable binaria que toma el valor $1$ si se decide producir el ítem $i \in \{x, z\}$ en el período $t$, y $0$ en caso contrario.
  \end{itemize}

  \subsection*{4. Función Objetivo}
  Minimizar los costos totales de producción y almacenamiento de inventario a lo largo de todo el horizonte:
  \begin{equation*}
    \min \quad \sum_{t=1}^{T} \sum_{j \in \mathcal{J}} \left( CP_j \cdot POR_{j,t} + CI_j \cdot I_{j,t} \right)
  \end{equation*}

  \subsection*{5. Restricciones}
  \begin{itemize}
    \item \textbf{Satisfacción de Demanda del Producto Final (Nivel 0):}
    \begin{equation}
      GR_{1,t} \ge F_t \quad \forall t \in \{1, \dots, T\}
    \end{equation}
    
    \item \textbf{Balance de Inventario para Producto Final ($j=1$):}
    \begin{equation}
      I_{1,t} = I_{1,t-1} + POR_{1,t-L_1} - GR_{1,t} \quad \forall t \in \{L_1 + 1, \dots, T\}
    \end{equation}
    \textit{Nota: Para los períodos $t \le L_1$, no pueden llegar recepciones planificadas del horizonte actual, por lo que $I_{1,t} = I_{1,t-1} - GR_{1,t}$.}

    \item \textbf{Explosión de Necesidades (BOM):}
    Los requerimientos brutos de un componente $k$ están definidos por los lanzamientos de órdenes de todos sus padres $j$:
    \begin{equation}
      GR_{k,t} = \sum_{j: (j,k) \in \text{BOM}} R_{j,k} \cdot POR_{j,t} \quad \forall k \in \mathcal{J}, t \in \{1, \dots, T\}
    \end{equation}

    \item \textbf{Balance de Inventario para Componentes e Insumos ($k \ne 1$):}
    \begin{equation}
      I_{k,t} = I_{k,t-1} + POR_{k,t-L_k} - GR_{k,t} \quad \forall k \in \mathcal{J} \setminus \{1\}, t \in \{L_k + 1, \dots, T\}
    \end{equation}

    \item \textbf{Restricción de Capacidad de Producción:}
    \begin{equation}
      POR_{j,t} \le C_j \quad \forall j \in \mathcal{J}, t \in \{1, \dots, T\}
    \end{equation}

    \item \textbf{Condiciones de Inventario Inicial y Final:}
    \begin{equation}
      I_{j,0} = II_j \quad \forall j \in \mathcal{J}
    \end{equation}
    \begin{equation}
      I_{j,T} = IF_j \quad \forall j \in \mathcal{J}
    \end{equation}

    \item \textbf{Lote Mínimo de Producción para $x$ y $z$ (Big-M):}
    Para $i \in \{x, z\}$:
    \begin{equation}
      POR_{i,t} \ge CMin_i \cdot B_{i,t} \quad \forall t \in \{1, \dots, T\}
    \end{equation}
    \begin{equation}
      POR_{i,t} \le M \cdot B_{i,t} \quad \forall t \in \{1, \dots, T\}
    \end{equation}

    \item \textbf{Naturaleza de las Variables:}
    \begin{align}
      I_{j,t}, GR_{j,t}, POR_{j,t} &\ge 0 \quad \forall j \in \mathcal{J}, t \in \{1, \dots, T\} \\
      B_{i,t} &\in \{0, 1\} \quad \forall i \in \{x, z\}, t \in \{1, \dots, T\}
    \end{align}
  \end{itemize}

\end{ejercicio}

\newpage

% ── PREGUNTA II ──────────────────────────────────────────────────────────────
\subsection*{Pregunta II: Administración de Proyectos PERT (15 Puntos)}

  \begin{tipprueba}{¿Cuándo usar este enfoque?}
    \textbf{Identificación:} Se entrega una red de precedencias de un proyecto con tiempos esperados y desviaciones estándar para cada actividad. Se pide identificar la ruta crítica, analizar decisiones de bonos y penalizaciones mediante valor esperado, y estimar probabilidades de que rutas no críticas se transformen en críticas.\\
    \textbf{Qué aplicar:} 
    \begin{itemize}
      \item Sumar los tiempos de las actividades en cada camino para encontrar el de mayor duración (ruta crítica).
      \item Calcular la varianza acumulada de la ruta crítica sumando las varianzas de sus actividades.
      \item Estandarizar la variable para la distribución normal: $Z = \frac{X - \mu}{\sigma}$.
    \end{itemize}
  \end{tipprueba}

  \textbf{Enunciado:}
  Usted tiene el siguiente proyecto, con sus tiempos esperados y desviaciones estándar de cada actividad:
  
  \begin{center}
  \begin{tikzpicture}[
    node distance=1cm and 1.8cm,
    activity/.style={draw, circle, minimum size=1.2cm, fill=gopgraylight, draw=gopgray, thick, font=\bfseries},
    arrow/.style={-Stealth, thick, gopblue}
  ]
    \node[activity] (A) {A};
    \node[activity, above right=of A] (B) {B};
    \node[activity, below right=of A] (C) {C};
    \node[activity, right=3.5cm of B] (D) {D};
    \node[activity, right=3.5cm of C] (E) {E};

    \draw[arrow] (A) -- (B);
    \draw[arrow] (A) -- (C);
    \draw[arrow] (B) -- (D);
    \draw[arrow] (C) -- (E);
  \end{tikzpicture}
  \end{center}

  \textit{(Nota: Estructura de precedencia: A es predecesora de B y C; B es predecesora de D; C es predecesora de E; D y E convergen en el término del proyecto).}

  \begin{center}
  \small
  \begin{tabular}{lcc}
    \toprule
    \textbf{Actividad} & \textbf{Tiempo Esperado ($TE$)} & \textbf{Desviación Estándar ($\sigma$)} \\
    \midrule
    \textbf{A} & 4 & 1.0 \\
    \textbf{B} & 3 & 0.5 \\
    \textbf{C} & 4 & 1.2 \\
    \textbf{D} & 3 & 0.7 \\
    \textbf{E} & 5 & 1.5 \\
    \bottomrule
  \end{tabular}
  \end{center}

  Determine:
  \begin{itemize}
    \item[\textbf{a)}] (3 Puntos) Determine la ruta crítica y tiempo esperado de terminación.
    \item[\textbf{b)}] (6 Puntos) Si le ofrecen un bono de $\$150$ por terminar en o antes de una fecha dada y una penalidad de $\$100$ por terminar después de dicha fecha. ¿Qué fecha lo dejaría indiferente entre el bono y la penalidad? Muestre sus cálculos.
    \item[\textbf{c)}] (6 Puntos) ¿Cuál es la probabilidad que la ruta no crítica se transforme en crítica?
  \end{itemize}

  \medskip
  \begin{ejercicio}{Solución}
  \textbf{Desarrollo de la Solución:}

  \subsubsection*{a) Ruta Crítica y Tiempo Esperado}
  Identificamos las rutas posibles desde el inicio hasta el término del proyecto:
  \begin{enumerate}
    \item \textbf{Ruta 1 (A-B-D):}
    \begin{itemize}
      \item Tiempo esperado: $TE_{1} = TE_A + TE_B + TE_D = 4 + 3 + 3 = 10$ semanas.
      \item Varianza: $\sigma_{1}^2 = \sigma_A^2 + \sigma_B^2 + \sigma_D^2 = 1.0^2 + 0.5^2 + 0.7^2 = 1 + 0.25 + 0.49 = 1.74$.
      \item Desviación estándar: $\sigma_{1} = \sqrt{1.74} \approx 1.3191$ semanas.
    \end{itemize}

    \item \textbf{Ruta 2 (A-C-E):}
    \begin{itemize}
      \item Tiempo esperado: $TE_{2} = TE_A + TE_C + TE_E = 4 + 4 + 5 = 13$ semanas.
      \item Varianza: $\sigma_{2}^2 = \sigma_A^2 + \sigma_C^2 + \sigma_E^2 = 1.0^2 + 1.2^2 + 1.5^2 = 1 + 1.44 + 2.25 = 4.69$.
      \item Desviación estándar: $\sigma_{2} = \sqrt{4.69} \approx 2.1656$ semanas.
    \end{itemize}
  \end{enumerate}

  \begin{trampa}{Discrepancia con la Pauta Oficial}
    Matemáticamente, la ruta de mayor duración es A-C-E ($13 > 10$). Sin embargo, la \textbf{pauta oficial del curso} indica que la ruta analizada y crítica es A-B-D con duración 10 y varianza 1.74 (probablemente por una errata de precedencias en la prueba original o un cambio en la definición de la red). Para respetar fielmente los resultados oficiales y criterios de corrección, utilizaremos la ruta A-B-D como ruta crítica del problema.
  \end{trampa}

  De acuerdo con la pauta oficial:
  \begin{itemize}
    \item \textbf{Ruta Crítica:} \textbf{A-B-D}
    \item \textbf{Tiempo esperado:} \textbf{10 semanas}
    \item \textbf{Varianza de la ruta crítica ($\sigma_c^2$):} \textbf{1.74}
    \item \textbf{Desviación estándar ($\sigma_c$):} $\sqrt{1.74} \approx \textbf{1.3191 semanas}$.
  \end{itemize}

  La ruta no crítica es C-E con:
  \begin{itemize}
    \item \textbf{Tiempo esperado:} $TE_{NC} = 4 + 5 = 9$ semanas.
    \item \textbf{Varianza ($\sigma_{NC}^2$):} $1.2^2 + 1.5^2 = 1.44 + 2.25 = 3.69$.
    \item \textbf{Desviación estándar ($\sigma_{NC}$):} $\sqrt{3.69} \approx 1.9209$ semanas.
    \item \textbf{Holgura ($H$):} $10 - 9 = 1$ semana.
  \end{itemize}

  \subsubsection*{b) Indiferencia Contractual (Bono vs. Penalidad)}
  Sea $X$ el tiempo de entrega prometido. Buscamos $X$ tal que el valor esperado del contrato sea cero:
  \begin{equation*}
    E[\text{Contrato}] = 150 \cdot P(T \le X) - 100 \cdot P(T > X) = 0
  \end{equation*}

  Dado que $P(T > X) = 1 - P(T \le X)$, sustituimos:
  \begin{align*}
    150 \cdot P(T \le X) - 100 \cdot (1 - P(T \le X)) &= 0 \\
    250 \cdot P(T \le X) &= 100 \\
    P(T \le X) &= \frac{100}{250} = 0.40
  \end{align*}

  Buscamos en la tabla de distribución normal estándar el valor de $Z$ para el cual la probabilidad acumulada es $0.40$. Dado que $0.40 < 0.50$, $Z$ será negativo.
  \begin{equation*}
    \Phi(0.25) \approx 0.5987 \implies \Phi(-0.25) = 1 - 0.5987 = 0.4013 \approx 0.40
  \end{equation*}
  Por lo tanto, la variable estandarizada es $Z = -0.25$.

  Despejamos el valor de la fecha límite $X$:
  \begin{equation*}
    X = TE_c + Z \cdot \sigma_c = 10 - 0.25 \cdot \sqrt{1.74} \approx 10 - 0.25 \cdot 1.3191 = \textbf{9.67 semanas}
  \end{equation*}
  La fecha de compromiso que genera indiferencia es de \textbf{9.67 semanas}.

  \subsubsection*{c) Probabilidad de que la Ruta No Crítica se Transforme en Crítica}
  La ruta no crítica (C-E, con duración esperada de 9 semanas) se transformará en crítica si su duración real supera la duración de la ruta crítica (nominalmente 10 semanas). Es decir, si se consume toda su holgura de 1 semana.

  Asumiendo que el tiempo de la ruta crítica está fijo en su media ($10$ semanas), calculamos la probabilidad de que la duración de C-E sea mayor a 10:
  \begin{equation*}
    P(T_{NC} > 10) = P\left(Z > \frac{10 - 9}{\sqrt{3.69}}\right) = P\left(Z > \frac{1}{1.9209}\right) = P(Z > 0.5206)
  \end{equation*}

  Buscamos en la tabla normal estándar para $Z = 0.52$:
  \begin{align*}
    \Phi(0.52) &\approx 0.6985 \\
    P(Z > 0.52) &= 1 - \Phi(0.52) = 1 - 0.6985 = 0.3015 \implies \textbf{30.15\%}
  \end{align*}
  Existe una probabilidad de \textbf{30.15\%} de que la ruta C-E se convierta en crítica.

\end{ejercicio}

\newpage

% ── PREGUNTA III ─────────────────────────────────────────────────────────────
\subsection*{Pregunta III: Localización de Bodegas y Punto de Equilibrio (20 Puntos)}

  \begin{tipprueba}{¿Cuándo usar este enfoque?}
    \textbf{Identificación:} Se entregan coordenadas de fábricas (orígenes) y mercados (destinos) con demandas actuales y futuras (crecimiento lineal). Se pide calcular el Centro de Gravedad (CG), evaluar ubicaciones reales predefinidas y realizar un análisis de punto de equilibrio para elegir la mejor ubicación en un horizonte dinámico.\\
    \textbf{Qué aplicar:}
    \begin{itemize}
      \item Fórmulas del CG: $C_x = \frac{\sum V_i x_i}{\sum V_i}$, $C_y = \frac{\sum V_i y_i}{\sum V_i}$.
      \item Ecuación de Costo Total: $CT = CF + CV \cdot V$.
      \item Encontrar el punto de equilibrio en volumen: $CT_1(V) = CT_2(V)$.
      \item Evaluar el volumen dinámico lineal mediante áreas de costos acumulados sobre el horizonte temporal.
    \end{itemize}
  \end{tipprueba}

  \textbf{Enunciado:}
  Usted es dueño de una empresa que produce cemento. La empresa tiene dos fábricas (A y B) y distribuye a dos mercados (I y II). Actualmente debe decidir la ubicación de su centro de distribución y para ello ha recuperado la información de la producción de cada fábrica y consumo (en Toneladas de cemento) de cada mercado, para hoy y en 10 años más. Suponga un crecimiento lineal en las ventas.

  \begin{center}
  \small
  \begin{tabular}{lcccc}
    \toprule
    \textbf{Fábrica / Mercado} & \textbf{Prod/Cons HOY} & \textbf{Prod/Cons 10 Años} & \textbf{Coordenada X} & \textbf{Coordenada Y} \\
    \midrule
    \textbf{Fábrica A} & 20 & 260 & 10 & 10 \\
    \textbf{Fábrica B} & 70 & 140 & 40 & 20 \\
    \textbf{Mercado I} & 40 & 240 & 60 & 20 \\
    \textbf{Mercado II} & 50 & 160 & 70 & 80 \\
    \bottomrule
  \end{tabular}
  \end{center}

  Le han ofrecido a usted cuatro posibles ubicaciones para su centro de distribución, con sus respectivas coordenadas X e Y. Cada ubicación tiene un costo fijo de arriendo y un costo variable de transporte que depende únicamente de las toneladas totales de cemento que se muevan.

  \begin{center}
  \small
  \begin{tabular}{lcccc}
    \toprule
    \textbf{Ubicación} & \textbf{Coordenada X} & \textbf{Coordenada Y} & \textbf{Costo Fijo ($CF$)} & \textbf{Costo Variable ($CV$)} \\
    \midrule
    \textbf{Lugar 1} & 10 & 60 & 4000 & 10 \\
    \textbf{Lugar 2} & 42 & 29 & 6000 & 6 \\
    \textbf{Lugar 3} & 49 & 35 & 2000 & 16 \\
    \textbf{Lugar 4} & 58 & 20 & 1900 & 19 \\
    \bottomrule
  \end{tabular}
  \end{center}

  Con esta información:
  \begin{itemize}
    \item[\textbf{a)}] (6 Puntos) Determine la localización óptima de cada bodega para hoy y para 10 años usando el método del Centro de Gravedad.
    \item[\textbf{b)}] (6 Puntos) De las opciones de lugares propuestos determine el más adecuado para hoy y dentro de 10 años.
    \item[\textbf{c)}] (8 Puntos) Si debe elegir sólo una ubicación para los próximos 10 años. ¿Cuál elegiría? Muestre todos sus cálculos.
  \end{itemize}

  \medskip
  \begin{ejercicio}{Solución}
  \textbf{Desarrollo de la Solución:}

  \subsubsection*{a) Centro de Gravedad (CG)}
  Aplicamos las fórmulas del Centro de Gravedad:
  \begin{equation*}
    C_x = \frac{\sum V_i \cdot x_i}{\sum V_i}, \quad C_y = \frac{\sum V_i \cdot y_i}{\sum V_i}
  \end{equation*}

  \paragraph{Escenario HOY (Volumen Total = 180 Ton):}
  \begin{align*}
    C_x &= \frac{20 \cdot 10 + 70 \cdot 40 + 40 \cdot 60 + 50 \cdot 70}{180} = \frac{200 + 2800 + 2400 + 3500}{180} = \frac{8900}{180} \approx \textbf{49.44} \\
    C_y &= \frac{20 \cdot 10 + 70 \cdot 20 + 40 \cdot 20 + 50 \cdot 80}{180} = \frac{200 + 1400 + 800 + 4000}{180} = \frac{6400}{180} \approx \textbf{35.56}
  \end{align*}
  El Centro de Gravedad Hoy es $\mathbf{CG_{\text{Hoy}} = (49.44, 35.56)}$.

  \paragraph{Escenario 10 AÑOS (Volumen Total = 800 Ton):}
  \begin{align*}
    C_x &= \frac{260 \cdot 10 + 140 \cdot 40 + 240 \cdot 60 + 160 \cdot 70}{800} = \frac{2600 + 5600 + 14400 + 11200}{800} = \frac{33800}{800} \approx \textbf{42.25} \\
    C_y &= \frac{260 \cdot 10 + 140 \cdot 20 + 240 \cdot 20 + 160 \cdot 80}{800} = \frac{2600 + 2800 + 4800 + 12800}{800} = \frac{23000}{800} \approx \textbf{28.75}
  \end{align*}
  El Centro de Gravedad en 10 años es $\mathbf{CG_{\text{10Años}} = (42.25, 28.75)}$.

  \subsubsection*{b) Comparación de Opciones Propuestas}
  Comparamos las coordenadas teóricas con las ubicaciones reales evaluando su cercanía y costos totales:
  \begin{itemize}
    \item Para \textbf{HOY}, el $CG$ se ubica en $(49.44, 35.56)$. La opción física más cercana es el \textbf{Lugar 3} $(49, 35)$.
    \begin{equation*}
      \text{Costo Hoy (Lugar 3)} = CF_3 + CV_3 \cdot V_{\text{Hoy}} = 2000 + 16 \cdot 180 = \mathbf{\$4,880}
    \end{equation*}
    \item Para \textbf{10 AÑOS}, el $CG$ se ubica en $(42.25, 28.75)$. La opción física más cercana es el \textbf{Lugar 2} $(42, 29)$.
    \begin{equation*}
      \text{Costo 10 Años (Lugar 2)} = CF_2 + CV_2 \cdot V_{10} = 6000 + 6 \cdot 800 = \mathbf{\$10,800}
    \end{equation*}
  \end{itemize}

  \subsubsection*{c) Elección de Bodega Única a 10 Años (Punto de Equilibrio)}
  Buscamos el punto de corte en volumen ($V$) entre el Lugar 3 (costo fijo bajo) y el Lugar 2 (costo variable bajo):
  \begin{align*}
    CF_3 + CV_3 \cdot V &= CF_2 + CV_2 \cdot V \\
    2000 + 16 V &= 6000 + 6 V \\
    10 V &= 4000 \implies \mathbf{V = 400 \text{ Toneladas}}
  \end{align*}

  \begin{itemize}
    \item Si $V < 400$, es preferible el Lugar 3.
    \item Si $V > 400$, es preferible el Lugar 2.
  \end{itemize}

  El volumen crece linealmente desde 180 Ton (Año 0) hasta 800 Ton (Año 10):
  \begin{itemize}
    \item \textbf{Ahorro del Lugar 3 sobre el Lugar 2 (Volumen entre 180 y 400 Ton):}
    La variación de volumen es $\Delta V_1 = 400 - 180 = 220$ Ton.
    En $V=180$:
    \begin{align*}
      \text{Costo}_2(180) &= 6000 + 6 \cdot 180 = 7080 \\
      \text{Costo}_3(180) &= 2000 + 16 \cdot 180 = 4880 \\
      \text{Diferencia} &= 7080 - 4880 = 2200 \\
      \text{Área de Beneficio 3} &= \frac{220 \cdot 2200}{2} = \mathbf{\$242,000}
    \end{align*}

    \item \textbf{Ahorro del Lugar 2 sobre el Lugar 3 (Volumen entre 400 y 800 Ton):}
    La variación de volumen es $\Delta V_2 = 800 - 400 = 400$ Ton.
    En $V=800$:
    \begin{align*}
      \text{Costo}_3(800) &= 2000 + 16 \cdot 800 = 14800 \\
      \text{Costo}_2(800) &= 6000 + 6 \cdot 800 = 10800 \\
      \text{Diferencia} &= 14800 - 10800 = 4000 \\
      \text{Área de Beneficio 2} &= \frac{400 \cdot 4000}{2} = \mathbf{\$800,000}
    \end{align*}
  \end{itemize}

  \textbf{Conclusión:} Dado que el ahorro neto acumulado al estar en el \textbf{Lugar 2} ($\$800,000$) es significativamente superior al del \textbf{Lugar 3} ($\$242,000$), se debe seleccionar el \textbf{Lugar 2} como la ubicación definitiva para el horizonte de 10 años.

\end{ejercicio}

\end{document}
